% Architecture -> H, sample -> interpolation subset; algorithm -> one selected element.
% Volume three architecture figures. Independent export only; not book fonts.
\colorlet{MutedBlueFill}{FigureBlueFill}
\colorlet{MutedGreenFill}{FigureGreenFill}
\colorlet{MutedPurple}{FigurePurple}
\colorlet{MutedPurpleFill}{FigurePurpleFill}
\tikzset{
 v3 picture/.style={x=1mm,y=1mm,font=\figurefont\fontsize{8}{10}\selectfont,text=FigureInk,line width=.65pt},
 v3 box/.style={draw=FigureBlue,fill=white,rounded corners=1.8pt,line width=.7pt,inner sep=2.5pt,align=center},
 v3 data/.style={v3 box,fill=FigureBlueFill},
 v3 model/.style={v3 box,draw=FigureGreen,fill=FigureGreenFill},
 v3 output/.style={v3 box,draw=FigureYellow,fill=FigureYellowFill},
 v3 risk/.style={v3 box,draw=FigureRed,fill=FigureRedFill},
 v3 group/.style={draw=FigureBlue,rounded corners=1.8pt,line width=.65pt,inner sep=3mm},
 v3 flow/.style={draw=FigureStroke,line width=.75pt,-{Stealth[length=1.65mm,width=1.05mm]},rounded corners=1pt},
 v3 guide/.style={draw=FigureMuted,line width=.6pt,dashed},
 v3 feedback/.style={v3 flow,draw=FigureGreen,dashed},
 v3 title/.style={draw=none,fill=none,font=\figurefont\fontsize{8}{10}\selectfont,align=center,inner sep=1pt},
 v3 note/.style={draw=none,fill=none,font=\FigureNoteFont\fontsize{7.5}{9.5}\selectfont,text=FigureMuted,align=center,inner sep=1pt},
 v3 label/.style={draw=none,fill=none,font=\figurefont\fontsize{8}{10}\selectfont,align=center,inner sep=1pt},
 v3 wire/.style={draw=FigureMuted,line width=.45pt},
 v3 point/.style={circle,draw=FigureBlue,fill=FigureBlueFill,minimum size=2.3mm,inner sep=0pt,line width=.55pt}
}
% Bring pre-existing native TikZ geometry to the same typography and scale.
\tikzset{
 foundation picture/.append style={font=\figurefont\fontsize{8}{10}\selectfont,line width=.65pt},
 foundation node/.append style={rounded corners=1.8pt,line width=.7pt,font=\figurefont\fontsize{8}{10}\selectfont},
 foundation heading/.append style={draw=none,fill=none,font=\figurefont\fontsize{8}{10}\selectfont},
 foundation note/.append style={draw=none,fill=none,font=\FigureNoteFont\fontsize{7.5}{9.5}\selectfont},
 foundation flow/.append style={line width=.75pt,-{Stealth[length=1.65mm,width=1.05mm]}},
 foundation frame/.append style={rounded corners=1.8pt,line width=.7pt},
 foundation edge label/.append style={draw=none,fill=none,font=\figurefont\fontsize{8}{10}\selectfont}
}
% Official Lucide outline icons, consistently used across diagrams. Size in mm.
\newcommand{\VIcon}[5]{\node[inner sep=0pt] at (#2,#3) {\includegraphics[width=#4mm]{figures/03-foundations/00-shared/tikz/icons/#1-#5.pdf}};}
% A neutral network glyph (not a claim about an actual model architecture).
% Coordinates denote its centre; scale 1 gives 13 x 8 mm.
\newcommand{\VNet}[4]{\begin{scope}[shift={(#1,#2)},scale=#3]
 \foreach \a in {-3,0,3}{\foreach \b in {-3,0,3}{\draw[v3 wire] (-5,\a)--(0,\b);}}
 \foreach \a in {-3,0,3}{\draw[v3 wire] (0,\a)--(5,0);}
 \foreach \a in {-3,0,3}{\draw[draw=#4,fill=#4,line width=.55pt] (-5,\a) circle(.9);\draw[draw=#4,fill=white,line width=.55pt] (0,\a) circle(.9);}
 \draw[draw=#4,fill=#4,line width=.55pt] (5,0) circle(.9);
\end{scope}}
\newcommand{\VDots}[4]{\foreach \k in {0,1,2,3}{\fill[#4] ({#1+2.7*\k},#2) circle(#3);}}
% Feature vector/tensor elements are scientific glyphs, not decorative icons.
\newcommand{\VVector}[4]{\foreach \k in {0,...,4}{\fill[#4,rounded corners=.25pt] ({#1+2.2*\k},#2) rectangle ++(1.5,#3);}}

\begin{tikzpicture}[v3 picture]
\path[use as bounding box] (0,0) rectangle (163,77);
\node[v3 model,minimum width=40mm,minimum height=16mm] (S) at (23,59) {};
\VIcon{image}{11}{59}{9}{blue}
\node[v3 label,anchor=west] at (20,59) {样本 $S$};
\foreach \y/\icon/\name in {40/dices/初始化,26/settings/优化器,12/shield-check/正则化}{
 \node[v3 output,minimum width=24mm,minimum height=10mm] (in\y) at (15,\y) {};
 \VIcon{\icon}{7}{\y}{5}{gold}\node[v3 label,anchor=west] at (12,\y) {\name};
}
\node[v3 data,minimum width=24mm,minimum height=16mm] (A) at (57,26) {};
\VNet{57}{26}{1}{FigureBlue}
\node[v3 label,anchor=north] at (57,16) {选解算法 $A$};
\draw[v3 flow] (in40.east)--(34,40)--(34,30)--(45,30);
\draw[v3 flow] (in26.east)--(A.west);
\draw[v3 flow] (in12.east)--(34,12)--(34,22)--(45,22);
\draw[v3 flow] (S.east)--(57,59)--(A.north);
% Sets are qualitative contours; neither area nor point count has quantitative meaning.
\path[draw=FigureBlue,fill=FigureBlueFill,line width=.8pt]
 (85,18) .. controls (71,32) and (80,49) .. (99,59)
 .. controls (109,68) and (116,73) .. (132,70)
 .. controls (151,70) and (160,62) .. (160,44)
 .. controls (164,25) and (155,11) .. (141,8)
 .. controls (115,3) and (97,7) .. (85,18)--cycle;
\node[v3 title] at (124,75) {架构限定假设空间 $\mathcal H$};
\path[draw=FigureBlue,fill=FigureGreenFill,dashed,line width=.75pt]
 (105,20) .. controls (103,27) and (94,35) .. (102,43)
 .. controls (109,50) and (116,44) .. (124,47)
 .. controls (136,53) and (142,45) .. (143,34)
 .. controls (147,24) and (139,17) .. (130,16)
 .. controls (119,13) and (109,15) .. (105,20)--cycle;
\foreach \x/\y in {88/33,93/25,94/44,101/53,108/60,116/66,125/61,136/64,145/58,152/51,155/40,150/27,145/19,138/12,125/10,112/11,101/16,107/24,110/39,118/35,127/42,136/36,139/27,130/21,119/20}{\fill[FigureBlue] (\x,\y) circle(.65);}
\draw[v3 flow] (A.east)--(116,26);
\fill[FigureGreen] (119,26) circle(1.7);\draw[FigureGreen,line width=.8pt] (119,26) circle(1.7);
\draw[v3 flow,dashed,draw=FigureGreen] (S.east)--(119,59)--(119,44);
\node[v3 note,fill=white,inner sep=1pt] at (77,55) {样本施加约束};
\node[v3 label] at (124,1.8) {插值解集合：$\hat R_S=0$};
\node[v3 label,anchor=west] at (143,4) {选中 $\hat h$};
\end{tikzpicture}
